% Propietario conservado: /Users/ruben/Documents/New project/output/EXTREMA_MEDIA_RAZON_RECIPROCIDAD_ES_EN_20260918/ARTICULO/ES/source/libro/certificacion_lean.tex
\chapter*{Arithmetic formalization of bilateral balance}
\addcontentsline{toc}{section}{Arithmetic formalization of bilateral balance}
\label{x_reciprocidad:libro:formalizacion-lean}

The nonadic specialization of the bilateral law admits an independent arithmetic
formalization. The point of application lies after the
construction of its coefficients and transports: the base-ten partition fixes
\(a=8\), and the two registers are \(8x-y\) and \(9x+y\). On that domain,
Lean~4.21.0 verifies seven theorems with its library \texttt{Std}. The formal
sources and execution procedure are provided as a supplement.
This complement states their mathematical content and relates it to
the operator proofs of the article.

\section*{Polynomial identity and conservation of the norm}

For any \(x,y\in\mathbb Z\), we have
\[
 9(8x-y)^2+8(9x+y)^2=17(72x^2+y^2).
\]
Expansion of the two squares cancels the cross products. When
\(y^2=x^2\), the same equality gives
\[
 9(8x-y)^2+8(9x+y)^2=17\cdot73\,x^2.
\]
These are the first two formalized theorems. The first expresses a
balance for arbitrary integer inputs; the second applies the explicit
hypothesis of norm conservation. The integer \(73=72+1\) occupies
exactly the place it has in \cref{x_reciprocidad:lib04:teo-balance}.

\section*{Arbitrary finite dimension}

Let \(n\in\mathbb N\) and \(x,y:\mathbb N\to\mathbb Z\). Define
\[
 S_n(x)=\sum_{j<n}x_j^2,
 \qquad
 E_n(x,y)=\sum_{j<n}\bigl(9(8x_j-y_j)^2+8(9x_j+y_j)^2\bigr).
\]
The next two theorems establish
\[
 E_n(x,y)=17\bigl(72S_n(x)+S_n(y)\bigr)
\]
and, under the hypothesis \(S_n(y)=S_n(x)\),
\[
 E_n(x,y)=17\cdot73\,S_n(x).
\]
The formal proof proceeds by induction on \(n\). The initial case is the
empty sum; the induction step adds a coordinate and applies the already proved
polynomial identity. The universal quantifier over \(n\) includes
all finite dimensions. The hypothesis uses the total sum of
squares, so it admits redistributions between coordinates.

\section*{Exact inversion of the bilateral channels}

Write \(q_-=8x-y\) and \(q_+=9x+y\). The fifth and sixth theorems
are the exact equalities
\[
 q_-+q_+=17x,\qquad 8q_+-9q_-=17y.
\]
In the image of the map \((x,y)\mapsto(q_-,q_+)\), both left-hand
sides are multiples of \(17\). The divisions recover the two
inputs. The seventh theorem states their injectivity consequence:
\[
 \left.
 \begin{aligned}
 8x-y&=8x'-y',\\
 9x+y&=9x'+y'
 \end{aligned}
 \right\}
 \quad\Longrightarrow\quad x=x'\ \text{and}\ y=y'.
\]
Cancellation in \(\mathbb Z\) justifies this conclusion. The formal
content coincides with recovery of the pair of registers, before
normalizing them as channels of a metric realization.

\section*{Correspondence between algebraic proof and formalization}

The formalization uses arbitrary integer variables and mathematical
induction. It therefore constitutes a proof of the preceding statements,
distinct from checking a finite set of examples. The system
accepts the proof terms without added axiomatic declarations or
deferred proofs. The supplement also records the logical
dependencies reported by the tool for each declaration.

The operator law in the body of the article has a broader domain:
isometries between Hilbert spaces, archives of arbitrary depth and
conjugated readers. Its proofs are found in
\cref{x_reciprocidad:lib04:teo-balance,x_reciprocidad:lib06:simetria,x_reciprocidad:lib06:profundidad}. The Lean proof
incorporated here certifies the stated arithmetic specialization, with
its quantifiers and initial structure. The provenance of the
coefficients, generation of the register and exceptional realizations
retain the preceding proofs on which they depend.
